%!TEX root = ../note.tex
% Source: ../chapter3_econ101_fall2026.pdf
\section{Supply and producer choice}
This chapter mirrors Chapter 2 with sellers in place of buyers.

\subsection{Individual supply}
\begin{defi}
  The \term{individual supply curve} of a seller $j$ is the function
  $p \mapsto q_j^S(p)$ giving the quantity supplied at each price $p$,
  ceteris paribus.
\end{defi}

\begin{eg}[Shell]\label{eg:shell}
  Shell's supply of gas (millions of litres per day):
  \[
    q^S_{\text{Shell}}(p) =
    \begin{cases}
      0 & p < 1.60,\\
      8 + 5(p - 1.60) & 1.60 \leq p \leq 2.40,
    \end{cases}
  \]
  so $q^S(1.60) = 8$, $q^S(1.80) = 9$, $q^S(2.00) = 10$, $q^S(2.20) = 11$ and
  $q^S(2.40) = 12$. Below \dollar{1.60} Shell stops supplying gas.
\end{eg}

\begin{law}[Law of supply]
  Quantity supplied rises when the price rises, and vice versa:
  \[
    p_1 < p_2 \implies q^S(p_1) \leq q^S(p_2).
  \]
  That is, supply curves are \emph{upward sloping}.
\end{law}

\subsection{Perfect competition}
\begin{defi}
  A market is \term{perfectly competitive} if
  \begin{enumerate}
    \item there are many buyers and sellers, each small relative to the
      market, and no barriers to entry or exit;
    \item all firms sell an identical product;
    \item there is perfect information;
    \item mobility costs are zero: buyers can switch sellers for free.
  \end{enumerate}
  An agent is a \term{price taker} if it treats the market price $p$ as
  given, as opposed to a \emph{price maker}.
\end{defi}

\begin{nprop}
  In a perfectly competitive market every firm sells at the market price $p$.
\end{nprop}
\begin{proof}[Argument]
  If a firm charges $p' > p$, then by (ii)--(iv) every buyer switches to a
  rival at no cost, so it sells nothing. Suppose instead that selling at some
  $p' < p$ were profitable. With perfect information and free entry, a
  competitor would already have taken that opportunity, so this cannot hold in
  a situation where nobody wants to change what they do. Hence $p$ is
  exogenous to each firm.
\end{proof}

\begin{remark}
  Perfectly competitive markets are rare in reality. The assumption is a
  simplification that makes the analysis tractable.
\end{remark}

\subsection{The rational rule for sellers}
\begin{defi}
  A firm's total cost splits as
  \[
    \TC(q) = \FC + \VC(q),
  \]
  where the \term{fixed cost} $\FC$ does not depend on output (e.g.\
  capital) and the \term{variable cost} $\VC(q)$ does (e.g.\ labour). The
  \term{marginal cost} is $\MC(q) = \TC(q) - \TC(q-1)$, or $\TC'(q)$ in the
  continuous case.
\end{defi}

\begin{nprop}
  $\MC(q) = \VC(q) - \VC(q-1)$. In particular, fixed costs do not affect the
  choice of quantity.
\end{nprop}
\begin{proof}
  $\FC$ cancels in the difference. This is Proposition~\ref{prop:sunk} again.
\end{proof}

For a price taker, the marginal benefit of selling one more unit is the price
$p$.

\begin{nthm}[Rational rule for sellers]\label{thm:seller}
  If $\MC$ is nondecreasing, the profit-maximising quantity at price $p$ is
  \[
    q^S(p) = \max\set{q \geq 0 : q = 0 \text{ or } \MC(q) \leq p}.
  \]
  At an interior optimum, $\text{price} = \text{marginal cost}$. So:
  \emph{sell one more unit if $p \geq \MC$.}
\end{nthm}
\begin{proof}
  Apply Theorem~\ref{thm:rational} with $\MB(q) = p$ constant.
\end{proof}

\begin{ncor}
  Increasing marginal cost implies the law of supply, and the supply curve
  \emph{is} the marginal cost curve.
\end{ncor}
\begin{proof}
  If $p_1 < p_2$ then $\set{q : \MC(q) \leq p_1} \subseteq \set{q : \MC(q)
  \leq p_2}$.
\end{proof}

\begin{remark}
  The law of supply comes from two sources: \emph{diminishing marginal
  product} (below) and \emph{rising input costs}.
\end{remark}

\subsection{Diminishing marginal product}
\begin{defi}
  A \term{production function} $Y = F(L, K)$ gives output from labour $L$
  (variable input) and capital $K$ (fixed input). The \term{marginal
  product of labour} is $\MP_L = \partial F / \partial L$, or
  $F(L, K) - F(L-1, K)$ in discrete terms. There is \term{diminishing
  marginal product} if $\MP_L$ decreases in $L$ for fixed $K$.
\end{defi}

\begin{nthm}[Diminishing marginal product $\Rightarrow$ rising marginal cost]
  \label{thm:mpmc}
  Let the wage be $w$ and $K$ fixed, and let $F(\cdot, K)$ be smooth and
  strictly increasing, with inverse $L(Y)$ (the labour needed to make $Y$).
  Then
  \[
    \MC(Y) = \frac{w}{\MP_L\bigl(L(Y)\bigr)}.
  \]
  Hence if $\MP_L$ is decreasing, $\MC$ is increasing.
\end{nthm}
\begin{proof}
  $\TC(Y) = wL(Y) + \FC$, so by the inverse function rule
  $\MC(Y) = wL'(Y) = w / F_L(L(Y), K)$. When $Y$ increases, $L(Y)$ increases,
  so $\MP_L$ falls and $\MC$ rises.
\end{proof}

\begin{significant}
  Each extra unit of output needs more and more labour, so each unit costs
  more than the last.
\end{significant}

\begin{eg}
  Take $Y = \sqrt{LK}$ with $K = 4$, so $Y = 2\sqrt{L}$ and
  $\MP_L = \frac12\sqrt{K/L} = 1/\sqrt{L}$, which is decreasing. In discrete
  terms:
  \begin{center}
    \begin{tabular}{cccccccc}
      \toprule
      $L$ & 0 & 1 & 2 & 3 & 4 & 5 & 6 \\
      $Y$ & 0 & 2 & 2.83 & 3.46 & 4.00 & 4.47 & 4.90 \\
      $\MP_L$ & --- & 2 & 0.83 & 0.64 & 0.54 & 0.47 & 0.43 \\
      \bottomrule
    \end{tabular}
  \end{center}
  Invert: $L = Y^2 / K = Y^2/4$. With $w = \dollar{15}$,
  $\TC(Y) = 15 \cdot Y^2/4 = 3.75\,Y^2$ (ignoring the fixed cost of capital),
  so
  \[
    \MC(Y) = \TC(Y) - \TC(Y-1) = 3.75\,(2Y - 1).
  \]
  \begin{center}
    \begin{tabular}{cccccccc}
      \toprule
      $Y$ & 0 & 1 & 2 & 3 & 4 & 5 & 6 \\
      $L$ & 0 & 0.25 & 1 & 2.25 & 4 & 6.25 & 9 \\
      $\TC$ & 0 & 3.75 & 15 & 33.75 & 60 & 93.75 & 135 \\
      $\MC$ & --- & 3.75 & 11.25 & 18.75 & 26.25 & 33.75 & 41.25 \\
      \bottomrule
    \end{tabular}
  \end{center}
  In the continuous version, $\MC(Y) = \TC'(Y) = 7.5\,Y$, which matches
  Theorem~\ref{thm:mpmc}: $w / \MP_L(L(Y)) = 15\sqrt{Y^2/4} = 7.5\,Y$. Setting
  $p = \MC$ gives the supply curve $q^S(p) = p / 7.5$, which is upward
  sloping.
\end{eg}
\begin{remark}
  The table's $L$ column assumes labour can be hired in fractions. $\MP_L$
  and $\MC$ move in opposite directions, as Theorem~\ref{thm:mpmc} predicts.
\end{remark}

\subsection{Market supply}
\begin{defi}
  The \term{market supply curve} is the horizontal sum
  $Q^S(p) = \sum_{j=1}^{M} q_j^S(p)$ over the $M$ sellers. If the sellers
  are identical, $Q^S(p) = M\,q^S(p)$.
\end{defi}

\begin{eg}
  With $M = 10$ identical firms:
  \begin{center}
    \begin{tabular}{cccccc}
      \toprule
      $p$ & 1.60 & 1.80 & 2.00 & 2.20 & 2.40 \\
      $q^S(p)$ & 6 & 8 & 10 & 12 & 14 \\
      $Q^S(p) = 10\,q^S(p)$ & 60 & 80 & 100 & 120 & 140 \\
      \bottomrule
    \end{tabular}
  \end{center}
  (millions of litres per day).
\end{eg}

A higher price raises market quantity supplied in two ways: existing firms
expand output, and new firms enter.

\subsection{Movements along versus shifts of supply}
As with demand, write
\[
  Q^S = Q^S(p;\ w,\ A,\ p_{\text{rel}},\ e,\ M),
\]
where $w$ is input prices, $A$ productivity, $p_{\text{rel}}$ prices of
related outputs, $e$ expectations and $M$ the number of sellers. A change in
$p$ is a \term{movement along the supply curve} and changes the
\term{quantity supplied}. A change in any other argument is a \term{shift
of the supply curve}. An increase in supply is a shift to the right.

\subsubsection*{The five supply shifters}
\begin{enumerate}
  \item \emph{Input prices.} Higher input prices raise $\MC$ at every $q$, so
    supply shifts left.
  \item \emph{Productivity and technology.} Producing more with the same
    inputs (new machinery, better management or processes, learning by doing)
    lowers $\MC$, so supply shifts right.
  \item \emph{Prices of related outputs.}
    \begin{defi}
      Good $y$ is a \term{substitute-in-production} for $x$ if it is made
      with the same inputs, and a \term{complement-in-production} if it is
      made together with $x$. Then
      \[
        \frac{\partial Q^S_x}{\partial p_y} < 0 \text{ for substitutes},\qquad
        \frac{\partial Q^S_x}{\partial p_y} > 0 \text{ for complements}.
      \]
    \end{defi}
  \item \emph{Expectations.} If firms expect the price to rise, they store
    output rather than sell it now, so supply today shifts left. This matters
    more for non-perishable goods (gasoline) than perishable ones (fresh
    fish).
  \item \emph{Type and number of sellers.} Changing $M$ shifts the market
    curve but no individual curve. For example, a Competition Bureau ruling on
    a merger changes $M$.
\end{enumerate}

\begin{eg}
  Crude oil is an input to gasoline. If its price rises, $\MC$ rises at every
  output level, so the market supply of gasoline shifts \emph{left}: at every
  price, less is supplied. On the slides, supply at \dollar{2.00} falls from
  $100$ to $60$ million litres.
\end{eg}

\subsection*{Summary: buyers versus sellers}
\begin{center}
  \begin{tabular}{lll}
    \toprule
    & Buyer & Seller \\
    \midrule
    Marginal benefit & $\MB(q)$ & $p$ \\
    Marginal cost & $p$ & $\MC(q)$ \\
    Rational rule & $p = \MB$ & $p = \MC$ \\
    Law & demand slopes down & supply slopes up \\
    Reason & diminishing $\MB$ & diminishing $\MP_L$, rising input costs \\
    Market curve & $\sum_i q_i^D$ & $\sum_j q_j^S$ \\
    \bottomrule
  \end{tabular}
\end{center}
